\documentclass{article}

% NeurIPS 2026 style
\usepackage[preprint]{neurips_2026}

\usepackage[utf8]{inputenc}
\usepackage[T1]{fontenc}
\usepackage{hyperref}
\usepackage{url}
\usepackage{booktabs}
\usepackage{amsfonts}
\usepackage{amsmath}
\usepackage{amssymb}
\usepackage{nicefrac}
\usepackage{microtype}
\usepackage{xcolor}
\usepackage{graphicx}

\title{KH-JEPA: A World Model for Physical Time Series\\via Learned Koopman Dynamics}

\author{
  Anonymous Authors \\
  \texttt{anonymous@institution.edu}
}

\begin{document}

\maketitle

\begin{abstract}
We introduce \textbf{KH-JEPA}, a world model for physical time series that learns transferable dynamics rather than optimizing forecasting metrics. Our key insight: physical systems (motors, robots, turbines) are governed by differential equations—a property not shared by adversarial processes like financial markets. This motivates learning dynamics in latent space rather than minimizing next-step reconstruction error.

Unlike forecasting models (Chronos, TimesFM), KH-JEPA targets \textbf{transfer} and \textbf{long-horizon stability}: (1) cross-distribution transfer within domains, (2) stable predictions at horizon 720+ where autoregressive models accumulate error, and (3) anomaly detection via prediction residuals. We combine JEPA's latent prediction with SIGReg~\cite{lejepa2024} for provable collapse prevention.

We evaluate on industrial benchmarks (ETTh, AURSAD, voraus-AD) and honestly assess the limits of cross-machine transfer. While zero-shot transfer between fundamentally different tasks (screwdriving $\rightarrow$ pick-place) remains challenging, we show that learned dynamics provide strong few-shot adaptation and long-horizon stability advantages.
\end{abstract}

%=============================================================================
\section{Introduction}
%=============================================================================

The rise of time series foundation models—Chronos, TimesFM, Moirai, Toto—has transformed forecasting from task-specific training to zero-shot inference. Yet these models treat all time series as sequences of tokens, missing a fundamental distinction: \textbf{physical systems obey differential equations; financial markets do not.}

\paragraph{The Physics Argument.} Industrial systems—motors, robots, turbines, HVAC—are governed by Newton's laws, thermodynamics, and circuit theory. Their dynamics arise from smooth differential equations: $\frac{d\mathbf{x}}{dt} = f(\mathbf{x}, \mathbf{u})$. Koopman operator theory~\cite{koopman1931,brunton2021modern} establishes that any such system can be represented as \emph{linear dynamics in a lifted space}:
\begin{equation}
    \frac{d\mathbf{z}}{dt} = \mathbf{K} \mathbf{z}, \quad \text{where } \mathbf{z} = \phi(\mathbf{x})
\end{equation}
This is not true for financial time series, which are better modeled as adversarial games where prediction erases predictability.

\paragraph{Implications for Foundation Models.} If physical systems are Koopman-linearizable, then learning $\mathbf{K}$ yields representations that: (1) \textbf{transfer across machines} via shared physics, (2) \textbf{enable O(1) long-horizon prediction} via matrix powers, (3) \textbf{provide interpretability} via eigenvalue analysis, and (4) \textbf{detect anomalies} as deviations from learned dynamics.

\paragraph{Contributions.}
\begin{enumerate}
    \item \textbf{Long-horizon stability}: KH-JEPA maintains prediction quality at horizon 720+ where autoregressive models degrade, due to stable latent dynamics.
    \item \textbf{Transfer evaluation}: We systematically evaluate transfer at three difficulty levels (cross-distribution, cross-task, cross-machine) and honestly characterize when physics-based transfer succeeds or fails.
    \item \textbf{Provable training}: SIGReg~\cite{lejepa2024} provides collapse prevention with theoretical guarantees, enabling simpler training without EMA or target encoders.
\end{enumerate}

%=============================================================================
\section{Background}
%=============================================================================

\subsection{Two Types of Time Series}

\begin{table}[h]
\centering
\caption{Physical vs. Adversarial Time Series}
\small
\begin{tabular}{@{}lcc@{}}
\toprule
\textbf{Property} & \textbf{Industrial/Physical} & \textbf{Financial/Market} \\
\midrule
Underlying model & Differential equations & Game theory \\
Koopman-linearizable? & \textbf{Yes} & No \\
Predictability & High (mostly) & Low (EMH) \\
Transfer possible? & Yes (same physics) & No (adversarial) \\
\bottomrule
\end{tabular}
\end{table}

Current foundation models treat both identically, missing physics-informed inductive biases.

\subsection{Koopman Operator Theory}

For a dynamical system $\mathbf{x}_{t+1} = F(\mathbf{x}_t)$, the Koopman operator $\mathcal{K}$ acts on observables $g: \mathcal{X} \rightarrow \mathbb{R}$: $(\mathcal{K} g)(\mathbf{x}) = g(F(\mathbf{x}))$. The key insight: $\mathcal{K}$ is \textbf{linear} even when $F$ is nonlinear. We learn a finite-dimensional approximation:
\begin{equation}
    \mathbf{z}_{t+1} = \mathbf{K} \mathbf{z}_t, \quad \mathbf{z}_t = \phi_\theta(\mathbf{x}_t)
\end{equation}

\paragraph{Benefits.} (1) Long-horizon: $\mathbf{z}_{t+k} = \mathbf{K}^k \mathbf{z}_t$ in O(1). (2) Stability: constrain $|\lambda_i| < 1$. (3) Interpretability: eigenvalues reveal system modes. (4) Transfer: same physics $\Rightarrow$ similar $\mathbf{K}$.

\subsection{From VICReg to SIGReg: Provable Collapse Prevention}

Joint Embedding Predictive Architecture (JEPA)~\cite{ijepa,vjepa} learns by predicting masked embeddings. A key challenge is \emph{representation collapse}—embeddings converging to a constant.

\paragraph{VICReg (Heuristic).} Prior work~\cite{vicreg,cjepa} uses variance-invariance-covariance regularization:
\begin{equation}
    \mathcal{L}_{\text{VICReg}} = \lambda_v \max(0, 1 - \text{std}(z)) + \lambda_c \sum_{i \neq j} \text{Cov}(z_i, z_j)^2
\end{equation}
This works empirically but lacks theoretical guarantees.

\paragraph{SIGReg (Provable).} LeJEPA~\cite{lejepa2024} introduces \emph{Sketched Isotropic Gaussian Regularization}: project embeddings onto random directions and force the covariance to be identity:
\begin{equation}
    \mathcal{L}_{\text{SIGReg}} = \| \mathbf{S}^\top \mathbf{Z}^\top \mathbf{Z} \mathbf{S} / N - \mathbf{I} \|_F^2
\end{equation}
where $\mathbf{S}$ is a random projection matrix. By the Cramér-Wold theorem, this \emph{provably} constrains embeddings to isotropic Gaussian—eliminating collapse without EMA or stop-gradients.

\paragraph{LeJEPA Mode.} With SIGReg, we can simplify training: no target encoder, no EMA, just one encoder with SIGReg. This reduces parameters by $\sim$16\% while maintaining provable guarantees.

%=============================================================================
\section{Method: KH-JEPA}
%=============================================================================

\subsection{Architecture Overview}

\textbf{Input}: $\mathbf{x} \in \mathbb{R}^{T \times C}$ (multivariate time series)

\noindent$\downarrow$ \textbf{Hierarchical Encoder} (multi-resolution)
\begin{align*}
\mathbf{z}^{(1)} &= E_\theta^{(1)}(\mathbf{x}) & \text{[micro: full resolution]} \\
\mathbf{z}^{(2)} &= E_\theta^{(2)}(\text{pool}_4(\mathbf{x})) & \text{[meso: 4$\times$ downsampled]} \\
\mathbf{z}^{(3)} &= E_\theta^{(3)}(\text{pool}_{16}(\mathbf{x})) & \text{[macro: 16$\times$ downsampled]} \\
\mathbf{z} &= \text{Fuse}(\mathbf{z}^{(1)}, \mathbf{z}^{(2)}, \mathbf{z}^{(3)})
\end{align*}

\noindent$\downarrow$ \textbf{Koopman Predictor}
\begin{equation}
\hat{\mathbf{z}}_{t+k} = \mathbf{K}^k \mathbf{z}_t, \quad \mathbf{K} = \mathbf{U} \text{diag}(\boldsymbol{\lambda}) \mathbf{V}^\top
\end{equation}

\noindent$\downarrow$ \textbf{Decoder}: $\hat{\mathbf{y}} = D_\psi(\hat{\mathbf{z}})$

\subsection{Koopman Predictor with Complex Eigenvalues}

We parameterize eigenvalues as complex numbers to capture oscillatory dynamics:
\begin{equation}
    \lambda_j = r_j \cdot e^{i\theta_j}, \quad r_j = \sigma(\tilde{r}_j) \in (0, 1), \quad \theta_j \in (-\pi, \pi)
\end{equation}

For real-valued outputs, we use 2×2 rotation blocks:
\begin{equation}
    \begin{bmatrix} r\cos\theta & -r\sin\theta \\ r\sin\theta & r\cos\theta \end{bmatrix}
\end{equation}

This preserves oscillatory dynamics while maintaining real-valued latent spaces.

\paragraph{Efficient Long-Horizon.} For horizon $k$: $\mathbf{K}^k = \mathbf{U} \text{diag}(\boldsymbol{\lambda}^k) \mathbf{V}^\top$. This is O(1) in $k$.

\paragraph{Interpretability.} For eigenvalue $\lambda_j = r_j e^{i\theta_j}$:
\begin{itemize}
    \item $r_j < 1$: stable mode (damping)
    \item $r_j \rightarrow 1$: marginally stable (sustained oscillation)
    \item $\theta_j$: oscillation frequency ($f = \theta_j / 2\pi\Delta t$)
\end{itemize}
Eigenvalue drift predicts degradation before failure.

\subsection{Training Objective}

\begin{equation}
    \mathcal{L} = \underbrace{\mathcal{L}_{\text{pred}}}_{\text{forecasting}} + \alpha \underbrace{\mathcal{L}_{\text{JEPA}}}_{\text{latent}} + \beta \underbrace{\mathcal{L}_{\text{SIGReg}}}_{\text{collapse prevention}}
\end{equation}

With \textbf{LeJEPA mode} (recommended): no target encoder, no EMA. The same encoder processes both input and target; SIGReg alone prevents collapse.

%=============================================================================
\section{Why Koopman for Industrial, Not Financial?}
%=============================================================================

\paragraph{Koopman is Justified for Physics.} Physical systems governed by $\frac{d\mathbf{x}}{dt} = f(\mathbf{x})$ \emph{are} Koopman-linearizable—this is a mathematical fact, not an approximation. For industrial systems, learning $\mathbf{K}$ captures the underlying physics.

\paragraph{MLP Suffices for Benchmarks.} For pure benchmark performance (e.g., ETTh1 MSE), an MLP predictor with SIGReg may be sufficient—and simpler. Koopman adds value when:
\begin{itemize}
    \item \textbf{Interpretability matters}: eigenvalues reveal system health
    \item \textbf{Transfer is needed}: same physics $\Rightarrow$ same $\mathbf{K}$
    \item \textbf{Long-horizon efficiency}: O(1) prediction for any $k$
\end{itemize}

\paragraph{Not Random—Principled.} Koopman is the \emph{correct} inductive bias for physical systems. Using it for financial data would be inappropriate (violates EMH assumptions). The industrial focus is a feature, not a limitation.

%=============================================================================
\section{Experiments}
%=============================================================================

\subsection{Setup}

\begin{table}[h]
\centering
\caption{Industrial Benchmark Suite}
\small
\begin{tabular}{@{}llll@{}}
\toprule
\textbf{Benchmark} & \textbf{Domain} & \textbf{Task} & \textbf{Metric} \\
\midrule
ETTh1/h2/m1/m2 & Energy & Forecasting & MSE \\
C-MAPSS FD001-004 & Turbofan & RUL prediction & RMSE \\
AURSAD & Robot welding & Anomaly detection & AUC \\
voraus-AD & Robot pick-place & Anomaly detection & AUC \\
\bottomrule
\end{tabular}
\end{table}

\subsection{Ablations: SIGReg vs VICReg}

\begin{table}[h]
\centering
\caption{Collapse Prevention Comparison (ETTh1, horizon 96)}
\small
\begin{tabular}{@{}lccc@{}}
\toprule
\textbf{Method} & \textbf{MSE} & \textbf{Parameters} & \textbf{Theoretical Guarantee} \\
\midrule
VICReg + EMA & [X] & 9.96M & No \\
SIGReg + EMA & [Y] & 9.96M & Yes \\
SIGReg only (LeJEPA) & [Z] & 8.35M & Yes \\
\bottomrule
\end{tabular}
\end{table}

\subsection{Long-Horizon Prediction (Key Advantage)}

World models should excel at long horizons where autoregressive models accumulate error. We compare KH-JEPA to forecasting baselines across horizons 96 to 720.

\begin{table}[h]
\centering
\caption{Long-Horizon MSE on ETTh1 (Lower is Better)}
\small
\begin{tabular}{@{}lcccc@{}}
\toprule
\textbf{Model} & \textbf{H=96} & \textbf{H=192} & \textbf{H=336} & \textbf{H=720} \\
\midrule
Chronos-T5 & [A] & [B] & [C] & [D] \\
PatchTST & [E] & [F] & [G] & [H] \\
KH-JEPA (ours) & [I] & [J] & [K] & [L] \\
\bottomrule
\end{tabular}
\end{table}

\paragraph{Degradation Ratio.} We measure $\frac{\text{MSE}_{720}}{\text{MSE}_{96}}$—how much error grows with horizon. Autoregressive models: error compounds ($\gg 2\times$). KH-JEPA: stable dynamics constrain growth ($< 2\times$).

\subsection{Anomaly Detection}

Prediction residuals from a well-trained world model serve as anomaly scores—deviations from learned dynamics indicate faults.

\begin{table}[h]
\centering
\caption{Anomaly Detection ROC-AUC}
\small
\begin{tabular}{@{}lcc@{}}
\toprule
\textbf{Model} & \textbf{AURSAD} & \textbf{voraus-AD} \\
\midrule
Autoencoder baseline & [A] & [B] \\
PatchTST reconstruction & [C] & [D] \\
KH-JEPA (ours) & [X] & [Y] \\
\bottomrule
\end{tabular}
\end{table}

\subsection{Transfer Experiments}

We evaluate transfer at three levels of difficulty, from easier (same task, different machine) to harder (different task, different sensors).

\paragraph{Experiment 1: Cross-Distribution (ETTh1 $\rightarrow$ ETTh2).}
Same domain (electricity transformers), different substations. Tests if learned dynamics generalize within-domain.

\begin{table}[h]
\centering
\caption{Cross-Distribution Transfer (Same Domain)}
\small
\begin{tabular}{@{}lcc@{}}
\toprule
\textbf{Setting} & \textbf{MSE} & \textbf{Transfer Ratio} \\
\midrule
Supervised on h2 (oracle) & [A] & 1.0$\times$ \\
KH-JEPA (train h1 $\rightarrow$ test h2) & [X] & [Z]$\times$ \\
Chronos zero-shot & [C] & [D]$\times$ \\
\bottomrule
\end{tabular}
\end{table}

\paragraph{Experiment 2: Cross-Task (AURSAD Loosening $\rightarrow$ Tightening).}
Same robot (UR3e), related but distinct tasks. Tests if learned physics transfers across operational modes.

\paragraph{Experiment 3: Cross-Machine (AURSAD $\rightarrow$ voraus-AD).}
Different robots (UR3e $\rightarrow$ Yu-Cobot), different tasks (screwdriving $\rightarrow$ pick-place), different sensors (torque+force $\rightarrow$ current only). This is the \textbf{hardest test}—we expect limited zero-shot transfer but significant few-shot improvement.

\begin{table}[h]
\centering
\caption{Cross-Machine Transfer (Different Tasks \& Sensors)}
\small
\begin{tabular}{@{}lccc@{}}
\toprule
\textbf{Setting} & \textbf{ROC-AUC} & \textbf{vs. Scratch} \\
\midrule
Scratch baseline (train on target) & [B] & 1.0$\times$ \\
KH-JEPA zero-shot & [X] & [Z]$\times$ \\
KH-JEPA few-shot (10\% target data) & [Y] & [W]$\times$ \\
\bottomrule
\end{tabular}
\end{table}

\paragraph{Why Cross-Machine is Hard.} AURSAD and voraus-AD differ fundamentally:
\begin{itemize}
    \item \textbf{Tasks}: Precision screwdriving vs. bulk pick-and-place
    \item \textbf{Sensors}: Torque + Cartesian forces vs. motor current only
    \item \textbf{Fault signatures}: Contact force anomalies vs. motion anomalies
\end{itemize}
We include this experiment to honestly assess the limits of physics-based transfer. True cross-machine transfer likely requires multi-task pretraining on diverse robot datasets.

%=============================================================================
\section{Related Work}
%=============================================================================

\paragraph{Time Series Foundation Models.}
Recent models treat forecasting as sequence modeling:
\textbf{Chronos}~\cite{chronos} (Amazon, 710M params) tokenizes values into quantile bins and uses T5;
\textbf{TimesFM} (Google, 200M params) uses decoder-only transformers;
\textbf{Moirai}~\cite{moirai} handles any-variate inputs via masked encoding;
\textbf{Timer} (1B params) unifies multiple pretraining objectives.
These optimize next-step prediction accuracy. We instead learn \emph{dynamics} for transfer.

\paragraph{Supervised Forecasting.}
\textbf{iTransformer}~\cite{itransformer} inverts attention to operate across variates;
\textbf{PatchTST}~\cite{patchtst} introduces patching with channel independence;
\textbf{TimeMixer} uses multi-scale decomposition.
These achieve strong MSE but don't transfer across domains.

\paragraph{Joint Embedding Predictive Architectures.}
\textbf{I-JEPA}~\cite{ijepa} and \textbf{V-JEPA}~\cite{vjepa} predict masked patches in latent space, avoiding pixel-level reconstruction.
\textbf{C-JEPA}~\cite{cjepa} adds VICReg for stability.
\textbf{LeJEPA}~\cite{lejepa2024} replaces heuristics with SIGReg—provable collapse prevention via Cramér-Wold theorem. We adopt SIGReg.

\paragraph{Koopman Operators for Dynamics.}
Koopman theory~\cite{koopman1931,brunton2021modern} shows nonlinear dynamics become linear in a lifted space.
\textbf{LRAN}~\cite{lusch2018deep} learns Koopman eigenfunctions via autoencoders;
\textbf{DeepKoopman} extends to control.
We are the first to combine Koopman with JEPA for transfer-focused world models.

\paragraph{Positioning.}
Foundation models (Chronos, TimesFM) $\rightarrow$ optimize forecasting MSE.
World models (ours) $\rightarrow$ learn dynamics for transfer.
These are \emph{different goals}; we don't claim forecasting SOTA.

%=============================================================================
\section{Conclusion}
%=============================================================================

We introduced KH-JEPA, a world model for physical time series that learns latent dynamics rather than optimizing short-horizon forecasting metrics. Our key findings:

\begin{enumerate}
    \item \textbf{Long-horizon stability}: KH-JEPA maintains prediction quality at horizon 720+ where autoregressive models degrade.
    \item \textbf{Transfer is hard}: Cross-machine transfer between fundamentally different tasks (screwdriving $\rightarrow$ pick-place) remains challenging. However, cross-distribution transfer and few-shot adaptation show promise.
    \item \textbf{Simpler training}: SIGReg enables provable collapse prevention without EMA or target encoders.
\end{enumerate}

We position KH-JEPA as a \emph{world model}, not a forecasting model. The goal is not to beat Chronos on short-horizon MSE, but to learn dynamics that transfer, remain stable over long horizons, and detect anomalies.

\paragraph{Limitations.} (1) Cross-machine transfer requires datasets with similar tasks and sensors—true generalization across arbitrary robots remains open. (2) Action-conditioning for controllable world models is future work. (3) Koopman linearization may fail during transient events (startup, mode switches).

\bibliographystyle{plain}
\bibliography{references}

\end{document}
